% =============================================================================
% FireSmokeGenNet --- Section IV (Results)
% Drop-in fragment for IEEEtran. Compile from the paper/ directory, or
% % =============================================================================
% FireSmokeGenNet --- Section IV (Results)
% Drop-in fragment for IEEEtran. Compile from the paper/ directory, or
% % =============================================================================
% FireSmokeGenNet --- Section IV (Results)
% Drop-in fragment for IEEEtran. Compile from the paper/ directory, or
% % =============================================================================
% FireSmokeGenNet --- Section IV (Results)
% Drop-in fragment for IEEEtran. Compile from the paper/ directory, or
% \input{paper/section_results} after setting \graphicspath{{paper/}}.
%
% SOURCE OF NUMBERS: results/json/ from
%   PYTHONPATH=src python3 -u scripts/run_pipeline.py --config configs/compact.yaml
% Regenerated tables: python scripts/export_results_tables.py
%
% This fragment REPLACES the previous Results section. Methods that were not
% executed in the compact public-data run are omitted (not copied from an
% earlier draft).
% =============================================================================

\section{Results}
\label{sec:results}

This section reports the compact public-data execution of FireSmokeGenNet.
It replaces the previous manuscript Results tables. All numerical values
are taken from the JSON measurements written by the compact pipeline.
Methods that were not run---Stable Diffusion, PowerPaint, BLD,
FlameDiffuser, official YOLO v6--v13 at \(640\times 640\), and registered
FLAME / HPWREN / SMOKE5K domain splits---are omitted rather than copied
from an earlier draft.

Images were \(64\times 64\) with a stride-4 TinyVAE (latent \(16\times 16\)).
The U-Net used widths 64--192, a cosine schedule with \(T=100\),
classifier-free guidance, and 8-step DDIM. The generator was trained for
160 iterations with MRDL \(\omega=0.4\). Quality ranking used the paper
weights \(Q=0.4\,\mathrm{color}+0.4\,\mathrm{visibility}+0.2\,\mathrm{translucency}\)
and retained the top 30\% of candidates. Source-level 80/10/10 splits used
camera identifiers as the source key, with SHA-256 exact-duplicate removal
and perceptual-hash near-duplicate control. Downstream inference used five
matched seeds \(\{42,123,3407,2025,2026\}\). Public substitutes were
Pyro-SDIS smoke frames and Wikimedia landscape backgrounds. Masks were
obtained by GrabCut inside the YOLO box followed by the largest connected
component. The host was CPU-only. This is a methodology-faithful run, not
a bit-identical replica of the \(512\)\,px / 20k-step / Qwen2-VL-7B /
YOLO@640 protocol.

Table~\ref{tab:protocol} records the realized sample counts.

\begin{table}[!t]
\centering
\caption{Compact public-data protocol actually executed.}
\label{tab:protocol}
\begin{tabular}{lr}
\hline
Quantity & Value \\
\hline
Real smoke train / val / test & 74 / 10 / 23 \\
Eligible backgrounds & 90 \\
Available masks & 74 \\
Generation pairs & 48 \\
Candidates & 96 \\
Retained (top 30\%) & 29 \\
\hline
\end{tabular}
\end{table}

\subsection{Vision--Language Quality-Predictor Validation}
\label{sec:results-vlm}

The original protocol specifies a LoRA-tuned Qwen2-VL-7B ranker on 150
human-rated images. That 7B teacher does not fit in the present CPU memory
budget. The same three axes, the same composite weights, the same top-30\%
keep rule, and the same five-fold student-versus-teacher protocol were
retained. The teacher was a heuristic scorer; the student was a small MLP
on image features.

Table~\ref{tab:vlm} reports mean \(\pm\) sample standard deviation across
the five folds. The student follows the teacher most closely on visibility
(MAE \(0.97\) on a 0--10 scale) and least closely on color (MAE \(2.71\)).
Composite rank correlation is weak and high-variance
(\(\rho=0.11\pm 0.25\); fold-wise
\(\rho\in\{0.45,-0.25,0.06,0.20,0.06\}\)).
These measurements support the pipeline structure---score three axes, form
\(Q\), keep the top 30\%---but do not reproduce the original Qwen2-VL-7B
agreement (\(\rho=0.83\)). The ranker is treated only as a within-study
filter on this 96-candidate pool. Fig.~\ref{fig:vlm} visualizes per-axis MAE.

\begin{table}[!t]
\centering
\caption{Five-fold cross-validation of the task-specific ranker.
Values are mean \(\pm\) sample SD across folds.
Spearman \(\rho\) is computed on the composite score \(Q\).}
\label{tab:vlm}
\begin{tabular}{lccc}
\hline
Quality dimension & MAE \(\downarrow\) & RMSE \(\downarrow\) & Spearman \(\rho\) \(\uparrow\) \\
\hline
Color fidelity & \(2.71\pm 0.47\) & \(3.64\pm 0.48\) & --- \\
Visibility & \(0.97\pm 0.18\) & \(1.26\pm 0.28\) & --- \\
Translucency & \(1.69\pm 0.26\) & \(2.16\pm 0.22\) & --- \\
Composite score & --- & --- & \(0.11\pm 0.25\) \\
\hline
\end{tabular}
\end{table}

\begin{figure}[!t]
\centering
\includegraphics[width=0.48\textwidth]{vlm_validation.png}
\caption{Five-fold MAE of the MLP student versus the three-axis heuristic
teacher. Error bars denote one sample standard deviation across folds.}
\label{fig:vlm}
\end{figure}

\subsection{Background Preservation and Text--Image Alignment}
\label{sec:results-bg}

PSNR, SSIM, and the 1-LPIPS proxy were computed over the preserved region
\(1-M\) after identical 64\,px resizing and \([0,1]\) normalization.
Prompt similarity is cosine similarity between the generated RGB image and
a bag-of-words embedding of the caption (CLIP ViT-H/14 was not loaded).
Only methods that were actually sampled are reported. Alpha-blend copies
the unmasked background exactly, so its PSNR/SSIM are not a fair
inpainting comparison.

\begin{table}[!t]
\centering
\caption{Background preservation and text--image alignment on the common
compact manifest. PSNR, SSIM, and the 1-LPIPS proxy are computed over
\(1-M\); prompt-sim is computed over the complete image.}
\label{tab:bg}
\begin{tabular}{lcccc}
\hline
Method & PSNR \(\uparrow\) & SSIM \(\uparrow\) & 1-LPIPS \(\uparrow\) & Prompt-sim \(\uparrow\) \\
\hline
Alpha-blend compositing & 60.00 & 1.00 & --- & --- \\
FireSmokeGenNet (raw decode) & 10.08 & 0.15 & 0.80 & 0.34 \\
\hline
\end{tabular}
\end{table}

The dual-branch generator, trained for 160 steps from a randomly
initialized TinyVAE, does not preserve the background at the original
paper-scale fidelity (PSNR \(28.10\)\,dB, SSIM \(0.88\)). The
\(10.08\)\,dB / \(0.15\) numbers measure this compact decoder, not smoke
realism inside the mask. They are reported because the metric definitions
match the experimental protocol; they are not a claim that the 64\,px
model matches SD-2 inpainting.

\subsection{Distribution Alignment}
\label{sec:results-dist}

Table~\ref{tab:dist} reports Fr\'echet distance and linear MMD\(^2\) in a
shared \(16\times 16\) RGB embedding of real test smoke versus synthetic
smoke. This is the compact analogue of CLIP ViT-H/14 CLIP-FD / linear MMD.
Lower is better. GAN, SD-inpainting, and FlameDiffuser rows are omitted
because those generators were not run.

\begin{table}[!t]
\centering
\caption{Feature-space distance between real Pyro-SDIS test smoke and
synthetic smoke. Lower is better. Features are \(16\times 16\) RGB, not
CLIP ViT-H/14.}
\label{tab:dist}
\begin{tabular}{lcc}
\hline
Method & Feature-FD \(\downarrow\) & MMD\(^2\) \(\downarrow\) \\
\hline
Alpha-blend compositing & 48.55 & 33.28 \\
FireSmokeGenNet & \textbf{44.54} & \textbf{28.83} \\
\hline
\end{tabular}
\end{table}

FireSmokeGenNet sits closer to real smoke than training-free compositing
on both distances (\(\Delta\) Feature-FD \(=-4.01\),
\(\Delta\) MMD\(^2=-4.45\)). The comparison is a single-manifest
descriptive result with no image-level confidence interval.
Fig.~\ref{fig:dist} visualizes the same values.

\begin{figure}[!t]
\centering
\includegraphics[width=0.48\textwidth]{distribution.png}
\caption{Feature-FD and linear MMD\(^2\) between real test smoke and
synthetic smoke. Lower is better.}
\label{fig:dist}
\end{figure}

\subsection{Boundary Softness Analysis}
\label{sec:results-bound}

The boundary band is
\(B_\delta(M)=\mathrm{Dilate}(M,\delta)\setminus\mathrm{Erode}(M,\delta)\)
with \(\delta\) scaled to 64\,px. Softness \(S\) is the mean gradient
magnitude in that band. KL uses 100-bin histograms of \(\|\nabla I\|_2\)
with a \(10^{-8}\) floor.

\begin{table}[!t]
\centering
\caption{Boundary-quality metrics. Real-smoke softness is a reference, not
a method.}
\label{tab:bound}
\begin{tabular}{lcc}
\hline
Method & Boundary softness \(S\) & KL\((\mathrm{real}\,\|\,\mathrm{synth})\) \(\downarrow\) \\
\hline
Real smoke & 0.363 & --- \\
Alpha-blend compositing & 0.438 & 15.20 \\
FireSmokeGenNet & \textbf{0.310} & 15.31 \\
\hline
\end{tabular}
\end{table}

FireSmokeGenNet is nearer the real reference than compositing
(\(0.310\) versus \(0.363\); blend \(0.438\)), i.e., less abrupt
mask-edge gradients. KL is comparable at this capacity (\(15.31\) versus
\(15.20\)). The originally published gap (softness \(0.124\) versus real
\(0.119\); KL \(0.041\) versus \(0.187\)) is not recovered at 64\,px /
160 steps. Softness is a local gradient descriptor, not a claim of
perceptual equivalence to real plumes. Fig.~\ref{fig:bound} shows the
three softness values.

\begin{figure}[!t]
\centering
\includegraphics[width=0.48\textwidth]{boundary.png}
\caption{Boundary-gradient magnitude in the \(\delta\)-band. Real smoke is
a reference distribution, not a generation method.}
\label{fig:bound}
\end{figure}

\subsection{Downstream Detection Performance}
\label{sec:results-det}

Official YOLO families at 640\,px were not executed. Pyro-SDIS median box
area is approximately \(0.17\%\) of the image; at 64--128\,px the box is
on the order of 10\,px, and IoU \(0.5\) is unstable. All eight compact
family stand-ins produced \(\mathrm{AP}_{50}=0\). Those zeros are not
filled into Table~\ref{tab:downstream}, and the previously published
YOLOv13 numbers (\(78.16\pm 0.27\%\rightarrow 80.88\pm 0.70\%\)) are not
copied.

The runnable downstream protocol is image-level smoke-versus-background
classification with a frozen ImageNet ResNet-18, five matched seeds, and a
test set of real smoke plus held-out backgrounds. Mixed training adds the
29 quality-filtered FireSmokeGenNet images to the real training set.

\begin{table*}[!t]
\centering
\caption{Five-seed comparison of real-only and mixed-data training.
Performance is smoke/background accuracy (\%). The upper endpoint of the
real-only \(t\)-interval exceeds 100\% because \(n=5\) and the sample SD
is large; the feasible range is \([82.22, 100]\).}
\label{tab:downstream}
\begin{tabular}{lccccc ccc ccc}
\hline
Condition & s1 & s2 & s3 & s4 & s5 & Mean \(\pm\) SD & 95\% CI & \(\Delta\) & Test & \(p\) & \(d_z\) \\
\hline
Real only
  & 81.40 & 86.05 & 97.67 & 100.00 & 100.00
  & \(93.02\pm 8.70\) & [82.22, 103.83]
  & --- & --- & --- & --- \\
Mixed data
  & 100.00 & 100.00 & 100.00 & 100.00 & 100.00
  & \(100.00\pm 0.00\) & [100.00, 100.00]
  & \(+6.98\) & Paired \(t\) & 0.147 & 0.80 \\
\hline
\end{tabular}
\end{table*}

The mixed-data mean is higher by 6.98 percentage points. Shapiro--Wilk on
the paired differences did not reject normality (\(p=0.111\)), so a
two-sided paired \(t\)-test was used. The test is not significant at
\(\alpha=0.05\) (\(p=0.147\)). Cohen's \(d_z=0.80\) is a large
standardized paired effect with only five runs; it is reported together
with the \(p\)-value and interval, not as a substitute for a detector AP
table. Holm correction is not applied because only one downstream
comparison was valid. Fig.~\ref{fig:cls} visualizes the two five-seed
means.

\begin{figure}[!t]
\centering
\includegraphics[width=0.48\textwidth]{classification.png}
\caption{Five-seed image-level smoke/background accuracy. Error bars
denote one sample standard deviation. Mixed training adds 29
quality-filtered FireSmokeGenNet images.}
\label{fig:cls}
\end{figure}

\subsection{Domain Generalization Under Environmental Shifts}
\label{sec:results-domain}

Table~\ref{tab:domain} is left empty. Summer\(\rightarrow\)winter,
day\(\rightarrow\)dusk, and clear\(\rightarrow\)haze splits on FLAME /
HPWREN / SMOKE5K were not downloaded because those sources require
registration. A compact clear-to-haze AP probe on the public subset also
returned degenerate \(\mathrm{AP}_{50}\) and is not reported as a
domain-shift result.

\begin{table}[!t]
\centering
\caption{Generalization under monitoring domain shifts. Not executed in
this public-data run.}
\label{tab:domain}
\begin{tabular}{lcccc}
\hline
Train \(\rightarrow\) Test & Real AP\(_{50}\) & Real Rec. & Mixed AP\(_{50}\) & Mixed Rec. \\
\hline
Summer \(\rightarrow\) Winter & --- & --- & --- & --- \\
Day \(\rightarrow\) Dusk & --- & --- & --- & --- \\
Clear \(\rightarrow\) Haze/Fog & --- & --- & --- & --- \\
\hline
\end{tabular}
\end{table}

\subsection{Qualitative Comparison and Ablations}
\label{sec:results-qual}

Fig.~\ref{fig:qual} compares FireSmokeGenNet with alpha-blend compositing
under matched backgrounds, masks, prompts, and DDIM settings. At 64\,px
the plumes occupy small regions, consistent with the Pyro-SDIS box-area
distribution. Where smoke is visible, the generator produces a different
interior appearance than constant-color blending rather than a copy of the
background.

\begin{figure*}[!t]
\centering
\includegraphics[width=0.95\textwidth]{qualitative.png}
\caption{Qualitative comparison under matched generation conditions. Each
row uses the same smoke-free background and target mask. Columns:
background, alpha-blend compositing, FireSmokeGenNet.}
\label{fig:qual}
\end{figure*}

JCA-off, MRDL-off, random encoder, unfiltered pool, and the full model
were trained under the compact detector protocol. All five configurations
produced \(\mathrm{AP}_{50}=0\). Table~\ref{tab:ablation} is therefore
unfilled. The modules remain implemented; they cannot be ranked by AP on
this host. The MRDL-weight sweep \(\omega\in\{0.0,0.4,1.0\}\) likewise
yielded \(\mathrm{AP}_{50}=0\) at every setting, so no AP curve is
reported. The training loss still uses
\((1-\omega)\mathcal{L}_{\mathrm{diff}}+\omega\mathcal{L}_{\mathrm{MRDL}}\)
with stop-gradient on the morphologically perturbed branch.

\begin{table}[!t]
\centering
\caption{Fixed-seed ablation of architectural components. AP columns are
left blank because the compact detector \(\mathrm{AP}_{50}\) was
degenerate.}
\label{tab:ablation}
\begin{tabular}{lcccc}
\hline
Configuration & AP\(_{50}\) & AP\(_{50:95}\) & Prec. & Rec. \\
\hline
No JCA (CA only) & --- & --- & --- & --- \\
No MRDL (\(\omega=0\)) & --- & --- & --- & --- \\
Encoder without ImageNet pretraining & --- & --- & --- & --- \\
Unfiltered candidate pool & --- & --- & --- & --- \\
Full FireSmokeGenNet & --- & --- & --- & --- \\
\hline
\end{tabular}
\end{table}

Table~\ref{tab:filter} reports candidate-pool sizes on the compact
percentage grid. AP and recall are unfilled. The 30\% operating point used
for generation retained 29 of 96 candidates (the configured 0.30 fraction
on the sampled pool).

\begin{table}[!t]
\centering
\caption{Retention-threshold sweep. Dataset size is measured; AP and
recall are not reported.}
\label{tab:filter}
\begin{tabular}{lccc}
\hline
Retention threshold & Dataset size & AP\(_{50}\) & Recall \\
\hline
10\% & 81 & --- & --- \\
20\% & 90 & --- & --- \\
30\% & 100 & --- & --- \\
40\% & 109 & --- & --- \\
50\% & 119 & --- & --- \\
No filtering & 167 & --- & --- \\
\hline
\end{tabular}
\end{table}

\subsection{Discussion}
\label{sec:results-disc}

\subsubsection{What this run supports}
By processing the mask and masked image through separate encoder
branches, dual-branch generation is closer to real Pyro-SDIS smoke than
alpha-blend compositing in Feature-FD and MMD\(^2\)
(Table~\ref{tab:dist}). Boundary softness is nearer the real reference
than compositing (Table~\ref{tab:bound}), while KL remains comparable at
this capacity. The three-axis composite and top-30\% keep rule were
executed; student--teacher \(\rho\) is too weak to claim a calibrated VLM
(Table~\ref{tab:vlm}). Mixed training raised five-seed classification
accuracy by \(+6.98\) percentage points (Table~\ref{tab:downstream}). The
paired test is not significant at \(\alpha=0.05\) on \(n=5\). This is a
directional proxy, not a YOLO \(\mathrm{AP}_{50}\) result.

\subsubsection{What this run does not support}
The compact execution does not support paper-scale background PSNR/SSIM
(\(28.10\)\,dB / \(0.88\)), YOLOv6--v13 \(\mathrm{AP}_{50}\) (including
the mixed YOLOv13 value of \(80.88\%\)), domain-shift AP on registered
wildfire-camera seasons, component / \(\omega\) / retention-threshold
rankings by AP, or Qwen2-VL-7B MAE/Spearman from the original VLM table.

\subsubsection{Limitations specific to the compact execution}
Tiny smoke boxes make IoU-0.5 detection degenerate. 160 generator steps
and a TinyVAE underfit background reconstruction. Heuristic quality
labels are not human ratings. The classification test set is small; mixed
accuracy saturates at 100\%, which compresses variance and inflates
\(d_z\) relative to a harder detector task.

\subsubsection{How to obtain the missing tables}
Run the paper-scale configuration on an A100-class GPU with official YOLO
weights, CLIP ViT-H/14, Qwen2-VL-7B, and the registered FLAME / HPWREN /
SMOKE5K sources. Until that run exists, those cells stay empty.
 after setting \graphicspath{{paper/}}.
%
% SOURCE OF NUMBERS: results/json/ from
%   PYTHONPATH=src python3 -u scripts/run_pipeline.py --config configs/compact.yaml
% Regenerated tables: python scripts/export_results_tables.py
%
% This fragment REPLACES the previous Results section. Methods that were not
% executed in the compact public-data run are omitted (not copied from an
% earlier draft).
% =============================================================================

\section{Results}
\label{sec:results}

This section reports the compact public-data execution of FireSmokeGenNet.
It replaces the previous manuscript Results tables. All numerical values
are taken from the JSON measurements written by the compact pipeline.
Methods that were not run---Stable Diffusion, PowerPaint, BLD,
FlameDiffuser, official YOLO v6--v13 at \(640\times 640\), and registered
FLAME / HPWREN / SMOKE5K domain splits---are omitted rather than copied
from an earlier draft.

Images were \(64\times 64\) with a stride-4 TinyVAE (latent \(16\times 16\)).
The U-Net used widths 64--192, a cosine schedule with \(T=100\),
classifier-free guidance, and 8-step DDIM. The generator was trained for
160 iterations with MRDL \(\omega=0.4\). Quality ranking used the paper
weights \(Q=0.4\,\mathrm{color}+0.4\,\mathrm{visibility}+0.2\,\mathrm{translucency}\)
and retained the top 30\% of candidates. Source-level 80/10/10 splits used
camera identifiers as the source key, with SHA-256 exact-duplicate removal
and perceptual-hash near-duplicate control. Downstream inference used five
matched seeds \(\{42,123,3407,2025,2026\}\). Public substitutes were
Pyro-SDIS smoke frames and Wikimedia landscape backgrounds. Masks were
obtained by GrabCut inside the YOLO box followed by the largest connected
component. The host was CPU-only. This is a methodology-faithful run, not
a bit-identical replica of the \(512\)\,px / 20k-step / Qwen2-VL-7B /
YOLO@640 protocol.

Table~\ref{tab:protocol} records the realized sample counts.

\begin{table}[!t]
\centering
\caption{Compact public-data protocol actually executed.}
\label{tab:protocol}
\begin{tabular}{lr}
\hline
Quantity & Value \\
\hline
Real smoke train / val / test & 74 / 10 / 23 \\
Eligible backgrounds & 90 \\
Available masks & 74 \\
Generation pairs & 48 \\
Candidates & 96 \\
Retained (top 30\%) & 29 \\
\hline
\end{tabular}
\end{table}

\subsection{Vision--Language Quality-Predictor Validation}
\label{sec:results-vlm}

The original protocol specifies a LoRA-tuned Qwen2-VL-7B ranker on 150
human-rated images. That 7B teacher does not fit in the present CPU memory
budget. The same three axes, the same composite weights, the same top-30\%
keep rule, and the same five-fold student-versus-teacher protocol were
retained. The teacher was a heuristic scorer; the student was a small MLP
on image features.

Table~\ref{tab:vlm} reports mean \(\pm\) sample standard deviation across
the five folds. The student follows the teacher most closely on visibility
(MAE \(0.97\) on a 0--10 scale) and least closely on color (MAE \(2.71\)).
Composite rank correlation is weak and high-variance
(\(\rho=0.11\pm 0.25\); fold-wise
\(\rho\in\{0.45,-0.25,0.06,0.20,0.06\}\)).
These measurements support the pipeline structure---score three axes, form
\(Q\), keep the top 30\%---but do not reproduce the original Qwen2-VL-7B
agreement (\(\rho=0.83\)). The ranker is treated only as a within-study
filter on this 96-candidate pool. Fig.~\ref{fig:vlm} visualizes per-axis MAE.

\begin{table}[!t]
\centering
\caption{Five-fold cross-validation of the task-specific ranker.
Values are mean \(\pm\) sample SD across folds.
Spearman \(\rho\) is computed on the composite score \(Q\).}
\label{tab:vlm}
\begin{tabular}{lccc}
\hline
Quality dimension & MAE \(\downarrow\) & RMSE \(\downarrow\) & Spearman \(\rho\) \(\uparrow\) \\
\hline
Color fidelity & \(2.71\pm 0.47\) & \(3.64\pm 0.48\) & --- \\
Visibility & \(0.97\pm 0.18\) & \(1.26\pm 0.28\) & --- \\
Translucency & \(1.69\pm 0.26\) & \(2.16\pm 0.22\) & --- \\
Composite score & --- & --- & \(0.11\pm 0.25\) \\
\hline
\end{tabular}
\end{table}

\begin{figure}[!t]
\centering
\includegraphics[width=0.48\textwidth]{vlm_validation.png}
\caption{Five-fold MAE of the MLP student versus the three-axis heuristic
teacher. Error bars denote one sample standard deviation across folds.}
\label{fig:vlm}
\end{figure}

\subsection{Background Preservation and Text--Image Alignment}
\label{sec:results-bg}

PSNR, SSIM, and the 1-LPIPS proxy were computed over the preserved region
\(1-M\) after identical 64\,px resizing and \([0,1]\) normalization.
Prompt similarity is cosine similarity between the generated RGB image and
a bag-of-words embedding of the caption (CLIP ViT-H/14 was not loaded).
Only methods that were actually sampled are reported. Alpha-blend copies
the unmasked background exactly, so its PSNR/SSIM are not a fair
inpainting comparison.

\begin{table}[!t]
\centering
\caption{Background preservation and text--image alignment on the common
compact manifest. PSNR, SSIM, and the 1-LPIPS proxy are computed over
\(1-M\); prompt-sim is computed over the complete image.}
\label{tab:bg}
\begin{tabular}{lcccc}
\hline
Method & PSNR \(\uparrow\) & SSIM \(\uparrow\) & 1-LPIPS \(\uparrow\) & Prompt-sim \(\uparrow\) \\
\hline
Alpha-blend compositing & 60.00 & 1.00 & --- & --- \\
FireSmokeGenNet (raw decode) & 10.08 & 0.15 & 0.80 & 0.34 \\
\hline
\end{tabular}
\end{table}

The dual-branch generator, trained for 160 steps from a randomly
initialized TinyVAE, does not preserve the background at the original
paper-scale fidelity (PSNR \(28.10\)\,dB, SSIM \(0.88\)). The
\(10.08\)\,dB / \(0.15\) numbers measure this compact decoder, not smoke
realism inside the mask. They are reported because the metric definitions
match the experimental protocol; they are not a claim that the 64\,px
model matches SD-2 inpainting.

\subsection{Distribution Alignment}
\label{sec:results-dist}

Table~\ref{tab:dist} reports Fr\'echet distance and linear MMD\(^2\) in a
shared \(16\times 16\) RGB embedding of real test smoke versus synthetic
smoke. This is the compact analogue of CLIP ViT-H/14 CLIP-FD / linear MMD.
Lower is better. GAN, SD-inpainting, and FlameDiffuser rows are omitted
because those generators were not run.

\begin{table}[!t]
\centering
\caption{Feature-space distance between real Pyro-SDIS test smoke and
synthetic smoke. Lower is better. Features are \(16\times 16\) RGB, not
CLIP ViT-H/14.}
\label{tab:dist}
\begin{tabular}{lcc}
\hline
Method & Feature-FD \(\downarrow\) & MMD\(^2\) \(\downarrow\) \\
\hline
Alpha-blend compositing & 48.55 & 33.28 \\
FireSmokeGenNet & \textbf{44.54} & \textbf{28.83} \\
\hline
\end{tabular}
\end{table}

FireSmokeGenNet sits closer to real smoke than training-free compositing
on both distances (\(\Delta\) Feature-FD \(=-4.01\),
\(\Delta\) MMD\(^2=-4.45\)). The comparison is a single-manifest
descriptive result with no image-level confidence interval.
Fig.~\ref{fig:dist} visualizes the same values.

\begin{figure}[!t]
\centering
\includegraphics[width=0.48\textwidth]{distribution.png}
\caption{Feature-FD and linear MMD\(^2\) between real test smoke and
synthetic smoke. Lower is better.}
\label{fig:dist}
\end{figure}

\subsection{Boundary Softness Analysis}
\label{sec:results-bound}

The boundary band is
\(B_\delta(M)=\mathrm{Dilate}(M,\delta)\setminus\mathrm{Erode}(M,\delta)\)
with \(\delta\) scaled to 64\,px. Softness \(S\) is the mean gradient
magnitude in that band. KL uses 100-bin histograms of \(\|\nabla I\|_2\)
with a \(10^{-8}\) floor.

\begin{table}[!t]
\centering
\caption{Boundary-quality metrics. Real-smoke softness is a reference, not
a method.}
\label{tab:bound}
\begin{tabular}{lcc}
\hline
Method & Boundary softness \(S\) & KL\((\mathrm{real}\,\|\,\mathrm{synth})\) \(\downarrow\) \\
\hline
Real smoke & 0.363 & --- \\
Alpha-blend compositing & 0.438 & 15.20 \\
FireSmokeGenNet & \textbf{0.310} & 15.31 \\
\hline
\end{tabular}
\end{table}

FireSmokeGenNet is nearer the real reference than compositing
(\(0.310\) versus \(0.363\); blend \(0.438\)), i.e., less abrupt
mask-edge gradients. KL is comparable at this capacity (\(15.31\) versus
\(15.20\)). The originally published gap (softness \(0.124\) versus real
\(0.119\); KL \(0.041\) versus \(0.187\)) is not recovered at 64\,px /
160 steps. Softness is a local gradient descriptor, not a claim of
perceptual equivalence to real plumes. Fig.~\ref{fig:bound} shows the
three softness values.

\begin{figure}[!t]
\centering
\includegraphics[width=0.48\textwidth]{boundary.png}
\caption{Boundary-gradient magnitude in the \(\delta\)-band. Real smoke is
a reference distribution, not a generation method.}
\label{fig:bound}
\end{figure}

\subsection{Downstream Detection Performance}
\label{sec:results-det}

Official YOLO families at 640\,px were not executed. Pyro-SDIS median box
area is approximately \(0.17\%\) of the image; at 64--128\,px the box is
on the order of 10\,px, and IoU \(0.5\) is unstable. All eight compact
family stand-ins produced \(\mathrm{AP}_{50}=0\). Those zeros are not
filled into Table~\ref{tab:downstream}, and the previously published
YOLOv13 numbers (\(78.16\pm 0.27\%\rightarrow 80.88\pm 0.70\%\)) are not
copied.

The runnable downstream protocol is image-level smoke-versus-background
classification with a frozen ImageNet ResNet-18, five matched seeds, and a
test set of real smoke plus held-out backgrounds. Mixed training adds the
29 quality-filtered FireSmokeGenNet images to the real training set.

\begin{table*}[!t]
\centering
\caption{Five-seed comparison of real-only and mixed-data training.
Performance is smoke/background accuracy (\%). The upper endpoint of the
real-only \(t\)-interval exceeds 100\% because \(n=5\) and the sample SD
is large; the feasible range is \([82.22, 100]\).}
\label{tab:downstream}
\begin{tabular}{lccccc ccc ccc}
\hline
Condition & s1 & s2 & s3 & s4 & s5 & Mean \(\pm\) SD & 95\% CI & \(\Delta\) & Test & \(p\) & \(d_z\) \\
\hline
Real only
  & 81.40 & 86.05 & 97.67 & 100.00 & 100.00
  & \(93.02\pm 8.70\) & [82.22, 103.83]
  & --- & --- & --- & --- \\
Mixed data
  & 100.00 & 100.00 & 100.00 & 100.00 & 100.00
  & \(100.00\pm 0.00\) & [100.00, 100.00]
  & \(+6.98\) & Paired \(t\) & 0.147 & 0.80 \\
\hline
\end{tabular}
\end{table*}

The mixed-data mean is higher by 6.98 percentage points. Shapiro--Wilk on
the paired differences did not reject normality (\(p=0.111\)), so a
two-sided paired \(t\)-test was used. The test is not significant at
\(\alpha=0.05\) (\(p=0.147\)). Cohen's \(d_z=0.80\) is a large
standardized paired effect with only five runs; it is reported together
with the \(p\)-value and interval, not as a substitute for a detector AP
table. Holm correction is not applied because only one downstream
comparison was valid. Fig.~\ref{fig:cls} visualizes the two five-seed
means.

\begin{figure}[!t]
\centering
\includegraphics[width=0.48\textwidth]{classification.png}
\caption{Five-seed image-level smoke/background accuracy. Error bars
denote one sample standard deviation. Mixed training adds 29
quality-filtered FireSmokeGenNet images.}
\label{fig:cls}
\end{figure}

\subsection{Domain Generalization Under Environmental Shifts}
\label{sec:results-domain}

Table~\ref{tab:domain} is left empty. Summer\(\rightarrow\)winter,
day\(\rightarrow\)dusk, and clear\(\rightarrow\)haze splits on FLAME /
HPWREN / SMOKE5K were not downloaded because those sources require
registration. A compact clear-to-haze AP probe on the public subset also
returned degenerate \(\mathrm{AP}_{50}\) and is not reported as a
domain-shift result.

\begin{table}[!t]
\centering
\caption{Generalization under monitoring domain shifts. Not executed in
this public-data run.}
\label{tab:domain}
\begin{tabular}{lcccc}
\hline
Train \(\rightarrow\) Test & Real AP\(_{50}\) & Real Rec. & Mixed AP\(_{50}\) & Mixed Rec. \\
\hline
Summer \(\rightarrow\) Winter & --- & --- & --- & --- \\
Day \(\rightarrow\) Dusk & --- & --- & --- & --- \\
Clear \(\rightarrow\) Haze/Fog & --- & --- & --- & --- \\
\hline
\end{tabular}
\end{table}

\subsection{Qualitative Comparison and Ablations}
\label{sec:results-qual}

Fig.~\ref{fig:qual} compares FireSmokeGenNet with alpha-blend compositing
under matched backgrounds, masks, prompts, and DDIM settings. At 64\,px
the plumes occupy small regions, consistent with the Pyro-SDIS box-area
distribution. Where smoke is visible, the generator produces a different
interior appearance than constant-color blending rather than a copy of the
background.

\begin{figure*}[!t]
\centering
\includegraphics[width=0.95\textwidth]{qualitative.png}
\caption{Qualitative comparison under matched generation conditions. Each
row uses the same smoke-free background and target mask. Columns:
background, alpha-blend compositing, FireSmokeGenNet.}
\label{fig:qual}
\end{figure*}

JCA-off, MRDL-off, random encoder, unfiltered pool, and the full model
were trained under the compact detector protocol. All five configurations
produced \(\mathrm{AP}_{50}=0\). Table~\ref{tab:ablation} is therefore
unfilled. The modules remain implemented; they cannot be ranked by AP on
this host. The MRDL-weight sweep \(\omega\in\{0.0,0.4,1.0\}\) likewise
yielded \(\mathrm{AP}_{50}=0\) at every setting, so no AP curve is
reported. The training loss still uses
\((1-\omega)\mathcal{L}_{\mathrm{diff}}+\omega\mathcal{L}_{\mathrm{MRDL}}\)
with stop-gradient on the morphologically perturbed branch.

\begin{table}[!t]
\centering
\caption{Fixed-seed ablation of architectural components. AP columns are
left blank because the compact detector \(\mathrm{AP}_{50}\) was
degenerate.}
\label{tab:ablation}
\begin{tabular}{lcccc}
\hline
Configuration & AP\(_{50}\) & AP\(_{50:95}\) & Prec. & Rec. \\
\hline
No JCA (CA only) & --- & --- & --- & --- \\
No MRDL (\(\omega=0\)) & --- & --- & --- & --- \\
Encoder without ImageNet pretraining & --- & --- & --- & --- \\
Unfiltered candidate pool & --- & --- & --- & --- \\
Full FireSmokeGenNet & --- & --- & --- & --- \\
\hline
\end{tabular}
\end{table}

Table~\ref{tab:filter} reports candidate-pool sizes on the compact
percentage grid. AP and recall are unfilled. The 30\% operating point used
for generation retained 29 of 96 candidates (the configured 0.30 fraction
on the sampled pool).

\begin{table}[!t]
\centering
\caption{Retention-threshold sweep. Dataset size is measured; AP and
recall are not reported.}
\label{tab:filter}
\begin{tabular}{lccc}
\hline
Retention threshold & Dataset size & AP\(_{50}\) & Recall \\
\hline
10\% & 81 & --- & --- \\
20\% & 90 & --- & --- \\
30\% & 100 & --- & --- \\
40\% & 109 & --- & --- \\
50\% & 119 & --- & --- \\
No filtering & 167 & --- & --- \\
\hline
\end{tabular}
\end{table}

\subsection{Discussion}
\label{sec:results-disc}

\subsubsection{What this run supports}
By processing the mask and masked image through separate encoder
branches, dual-branch generation is closer to real Pyro-SDIS smoke than
alpha-blend compositing in Feature-FD and MMD\(^2\)
(Table~\ref{tab:dist}). Boundary softness is nearer the real reference
than compositing (Table~\ref{tab:bound}), while KL remains comparable at
this capacity. The three-axis composite and top-30\% keep rule were
executed; student--teacher \(\rho\) is too weak to claim a calibrated VLM
(Table~\ref{tab:vlm}). Mixed training raised five-seed classification
accuracy by \(+6.98\) percentage points (Table~\ref{tab:downstream}). The
paired test is not significant at \(\alpha=0.05\) on \(n=5\). This is a
directional proxy, not a YOLO \(\mathrm{AP}_{50}\) result.

\subsubsection{What this run does not support}
The compact execution does not support paper-scale background PSNR/SSIM
(\(28.10\)\,dB / \(0.88\)), YOLOv6--v13 \(\mathrm{AP}_{50}\) (including
the mixed YOLOv13 value of \(80.88\%\)), domain-shift AP on registered
wildfire-camera seasons, component / \(\omega\) / retention-threshold
rankings by AP, or Qwen2-VL-7B MAE/Spearman from the original VLM table.

\subsubsection{Limitations specific to the compact execution}
Tiny smoke boxes make IoU-0.5 detection degenerate. 160 generator steps
and a TinyVAE underfit background reconstruction. Heuristic quality
labels are not human ratings. The classification test set is small; mixed
accuracy saturates at 100\%, which compresses variance and inflates
\(d_z\) relative to a harder detector task.

\subsubsection{How to obtain the missing tables}
Run the paper-scale configuration on an A100-class GPU with official YOLO
weights, CLIP ViT-H/14, Qwen2-VL-7B, and the registered FLAME / HPWREN /
SMOKE5K sources. Until that run exists, those cells stay empty.
 after setting \graphicspath{{paper/}}.
%
% SOURCE OF NUMBERS: results/json/ from
%   PYTHONPATH=src python3 -u scripts/run_pipeline.py --config configs/compact.yaml
% Regenerated tables: python scripts/export_results_tables.py
%
% This fragment REPLACES the previous Results section. Methods that were not
% executed in the compact public-data run are omitted (not copied from an
% earlier draft).
% =============================================================================

\section{Results}
\label{sec:results}

This section reports the compact public-data execution of FireSmokeGenNet.
It replaces the previous manuscript Results tables. All numerical values
are taken from the JSON measurements written by the compact pipeline.
Methods that were not run---Stable Diffusion, PowerPaint, BLD,
FlameDiffuser, official YOLO v6--v13 at \(640\times 640\), and registered
FLAME / HPWREN / SMOKE5K domain splits---are omitted rather than copied
from an earlier draft.

Images were \(64\times 64\) with a stride-4 TinyVAE (latent \(16\times 16\)).
The U-Net used widths 64--192, a cosine schedule with \(T=100\),
classifier-free guidance, and 8-step DDIM. The generator was trained for
160 iterations with MRDL \(\omega=0.4\). Quality ranking used the paper
weights \(Q=0.4\,\mathrm{color}+0.4\,\mathrm{visibility}+0.2\,\mathrm{translucency}\)
and retained the top 30\% of candidates. Source-level 80/10/10 splits used
camera identifiers as the source key, with SHA-256 exact-duplicate removal
and perceptual-hash near-duplicate control. Downstream inference used five
matched seeds \(\{42,123,3407,2025,2026\}\). Public substitutes were
Pyro-SDIS smoke frames and Wikimedia landscape backgrounds. Masks were
obtained by GrabCut inside the YOLO box followed by the largest connected
component. The host was CPU-only. This is a methodology-faithful run, not
a bit-identical replica of the \(512\)\,px / 20k-step / Qwen2-VL-7B /
YOLO@640 protocol.

Table~\ref{tab:protocol} records the realized sample counts.

\begin{table}[!t]
\centering
\caption{Compact public-data protocol actually executed.}
\label{tab:protocol}
\begin{tabular}{lr}
\hline
Quantity & Value \\
\hline
Real smoke train / val / test & 74 / 10 / 23 \\
Eligible backgrounds & 90 \\
Available masks & 74 \\
Generation pairs & 48 \\
Candidates & 96 \\
Retained (top 30\%) & 29 \\
\hline
\end{tabular}
\end{table}

\subsection{Vision--Language Quality-Predictor Validation}
\label{sec:results-vlm}

The original protocol specifies a LoRA-tuned Qwen2-VL-7B ranker on 150
human-rated images. That 7B teacher does not fit in the present CPU memory
budget. The same three axes, the same composite weights, the same top-30\%
keep rule, and the same five-fold student-versus-teacher protocol were
retained. The teacher was a heuristic scorer; the student was a small MLP
on image features.

Table~\ref{tab:vlm} reports mean \(\pm\) sample standard deviation across
the five folds. The student follows the teacher most closely on visibility
(MAE \(0.97\) on a 0--10 scale) and least closely on color (MAE \(2.71\)).
Composite rank correlation is weak and high-variance
(\(\rho=0.11\pm 0.25\); fold-wise
\(\rho\in\{0.45,-0.25,0.06,0.20,0.06\}\)).
These measurements support the pipeline structure---score three axes, form
\(Q\), keep the top 30\%---but do not reproduce the original Qwen2-VL-7B
agreement (\(\rho=0.83\)). The ranker is treated only as a within-study
filter on this 96-candidate pool. Fig.~\ref{fig:vlm} visualizes per-axis MAE.

\begin{table}[!t]
\centering
\caption{Five-fold cross-validation of the task-specific ranker.
Values are mean \(\pm\) sample SD across folds.
Spearman \(\rho\) is computed on the composite score \(Q\).}
\label{tab:vlm}
\begin{tabular}{lccc}
\hline
Quality dimension & MAE \(\downarrow\) & RMSE \(\downarrow\) & Spearman \(\rho\) \(\uparrow\) \\
\hline
Color fidelity & \(2.71\pm 0.47\) & \(3.64\pm 0.48\) & --- \\
Visibility & \(0.97\pm 0.18\) & \(1.26\pm 0.28\) & --- \\
Translucency & \(1.69\pm 0.26\) & \(2.16\pm 0.22\) & --- \\
Composite score & --- & --- & \(0.11\pm 0.25\) \\
\hline
\end{tabular}
\end{table}

\begin{figure}[!t]
\centering
\includegraphics[width=0.48\textwidth]{vlm_validation.png}
\caption{Five-fold MAE of the MLP student versus the three-axis heuristic
teacher. Error bars denote one sample standard deviation across folds.}
\label{fig:vlm}
\end{figure}

\subsection{Background Preservation and Text--Image Alignment}
\label{sec:results-bg}

PSNR, SSIM, and the 1-LPIPS proxy were computed over the preserved region
\(1-M\) after identical 64\,px resizing and \([0,1]\) normalization.
Prompt similarity is cosine similarity between the generated RGB image and
a bag-of-words embedding of the caption (CLIP ViT-H/14 was not loaded).
Only methods that were actually sampled are reported. Alpha-blend copies
the unmasked background exactly, so its PSNR/SSIM are not a fair
inpainting comparison.

\begin{table}[!t]
\centering
\caption{Background preservation and text--image alignment on the common
compact manifest. PSNR, SSIM, and the 1-LPIPS proxy are computed over
\(1-M\); prompt-sim is computed over the complete image.}
\label{tab:bg}
\begin{tabular}{lcccc}
\hline
Method & PSNR \(\uparrow\) & SSIM \(\uparrow\) & 1-LPIPS \(\uparrow\) & Prompt-sim \(\uparrow\) \\
\hline
Alpha-blend compositing & 60.00 & 1.00 & --- & --- \\
FireSmokeGenNet (raw decode) & 10.08 & 0.15 & 0.80 & 0.34 \\
\hline
\end{tabular}
\end{table}

The dual-branch generator, trained for 160 steps from a randomly
initialized TinyVAE, does not preserve the background at the original
paper-scale fidelity (PSNR \(28.10\)\,dB, SSIM \(0.88\)). The
\(10.08\)\,dB / \(0.15\) numbers measure this compact decoder, not smoke
realism inside the mask. They are reported because the metric definitions
match the experimental protocol; they are not a claim that the 64\,px
model matches SD-2 inpainting.

\subsection{Distribution Alignment}
\label{sec:results-dist}

Table~\ref{tab:dist} reports Fr\'echet distance and linear MMD\(^2\) in a
shared \(16\times 16\) RGB embedding of real test smoke versus synthetic
smoke. This is the compact analogue of CLIP ViT-H/14 CLIP-FD / linear MMD.
Lower is better. GAN, SD-inpainting, and FlameDiffuser rows are omitted
because those generators were not run.

\begin{table}[!t]
\centering
\caption{Feature-space distance between real Pyro-SDIS test smoke and
synthetic smoke. Lower is better. Features are \(16\times 16\) RGB, not
CLIP ViT-H/14.}
\label{tab:dist}
\begin{tabular}{lcc}
\hline
Method & Feature-FD \(\downarrow\) & MMD\(^2\) \(\downarrow\) \\
\hline
Alpha-blend compositing & 48.55 & 33.28 \\
FireSmokeGenNet & \textbf{44.54} & \textbf{28.83} \\
\hline
\end{tabular}
\end{table}

FireSmokeGenNet sits closer to real smoke than training-free compositing
on both distances (\(\Delta\) Feature-FD \(=-4.01\),
\(\Delta\) MMD\(^2=-4.45\)). The comparison is a single-manifest
descriptive result with no image-level confidence interval.
Fig.~\ref{fig:dist} visualizes the same values.

\begin{figure}[!t]
\centering
\includegraphics[width=0.48\textwidth]{distribution.png}
\caption{Feature-FD and linear MMD\(^2\) between real test smoke and
synthetic smoke. Lower is better.}
\label{fig:dist}
\end{figure}

\subsection{Boundary Softness Analysis}
\label{sec:results-bound}

The boundary band is
\(B_\delta(M)=\mathrm{Dilate}(M,\delta)\setminus\mathrm{Erode}(M,\delta)\)
with \(\delta\) scaled to 64\,px. Softness \(S\) is the mean gradient
magnitude in that band. KL uses 100-bin histograms of \(\|\nabla I\|_2\)
with a \(10^{-8}\) floor.

\begin{table}[!t]
\centering
\caption{Boundary-quality metrics. Real-smoke softness is a reference, not
a method.}
\label{tab:bound}
\begin{tabular}{lcc}
\hline
Method & Boundary softness \(S\) & KL\((\mathrm{real}\,\|\,\mathrm{synth})\) \(\downarrow\) \\
\hline
Real smoke & 0.363 & --- \\
Alpha-blend compositing & 0.438 & 15.20 \\
FireSmokeGenNet & \textbf{0.310} & 15.31 \\
\hline
\end{tabular}
\end{table}

FireSmokeGenNet is nearer the real reference than compositing
(\(0.310\) versus \(0.363\); blend \(0.438\)), i.e., less abrupt
mask-edge gradients. KL is comparable at this capacity (\(15.31\) versus
\(15.20\)). The originally published gap (softness \(0.124\) versus real
\(0.119\); KL \(0.041\) versus \(0.187\)) is not recovered at 64\,px /
160 steps. Softness is a local gradient descriptor, not a claim of
perceptual equivalence to real plumes. Fig.~\ref{fig:bound} shows the
three softness values.

\begin{figure}[!t]
\centering
\includegraphics[width=0.48\textwidth]{boundary.png}
\caption{Boundary-gradient magnitude in the \(\delta\)-band. Real smoke is
a reference distribution, not a generation method.}
\label{fig:bound}
\end{figure}

\subsection{Downstream Detection Performance}
\label{sec:results-det}

Official YOLO families at 640\,px were not executed. Pyro-SDIS median box
area is approximately \(0.17\%\) of the image; at 64--128\,px the box is
on the order of 10\,px, and IoU \(0.5\) is unstable. All eight compact
family stand-ins produced \(\mathrm{AP}_{50}=0\). Those zeros are not
filled into Table~\ref{tab:downstream}, and the previously published
YOLOv13 numbers (\(78.16\pm 0.27\%\rightarrow 80.88\pm 0.70\%\)) are not
copied.

The runnable downstream protocol is image-level smoke-versus-background
classification with a frozen ImageNet ResNet-18, five matched seeds, and a
test set of real smoke plus held-out backgrounds. Mixed training adds the
29 quality-filtered FireSmokeGenNet images to the real training set.

\begin{table*}[!t]
\centering
\caption{Five-seed comparison of real-only and mixed-data training.
Performance is smoke/background accuracy (\%). The upper endpoint of the
real-only \(t\)-interval exceeds 100\% because \(n=5\) and the sample SD
is large; the feasible range is \([82.22, 100]\).}
\label{tab:downstream}
\begin{tabular}{lccccc ccc ccc}
\hline
Condition & s1 & s2 & s3 & s4 & s5 & Mean \(\pm\) SD & 95\% CI & \(\Delta\) & Test & \(p\) & \(d_z\) \\
\hline
Real only
  & 81.40 & 86.05 & 97.67 & 100.00 & 100.00
  & \(93.02\pm 8.70\) & [82.22, 103.83]
  & --- & --- & --- & --- \\
Mixed data
  & 100.00 & 100.00 & 100.00 & 100.00 & 100.00
  & \(100.00\pm 0.00\) & [100.00, 100.00]
  & \(+6.98\) & Paired \(t\) & 0.147 & 0.80 \\
\hline
\end{tabular}
\end{table*}

The mixed-data mean is higher by 6.98 percentage points. Shapiro--Wilk on
the paired differences did not reject normality (\(p=0.111\)), so a
two-sided paired \(t\)-test was used. The test is not significant at
\(\alpha=0.05\) (\(p=0.147\)). Cohen's \(d_z=0.80\) is a large
standardized paired effect with only five runs; it is reported together
with the \(p\)-value and interval, not as a substitute for a detector AP
table. Holm correction is not applied because only one downstream
comparison was valid. Fig.~\ref{fig:cls} visualizes the two five-seed
means.

\begin{figure}[!t]
\centering
\includegraphics[width=0.48\textwidth]{classification.png}
\caption{Five-seed image-level smoke/background accuracy. Error bars
denote one sample standard deviation. Mixed training adds 29
quality-filtered FireSmokeGenNet images.}
\label{fig:cls}
\end{figure}

\subsection{Domain Generalization Under Environmental Shifts}
\label{sec:results-domain}

Table~\ref{tab:domain} is left empty. Summer\(\rightarrow\)winter,
day\(\rightarrow\)dusk, and clear\(\rightarrow\)haze splits on FLAME /
HPWREN / SMOKE5K were not downloaded because those sources require
registration. A compact clear-to-haze AP probe on the public subset also
returned degenerate \(\mathrm{AP}_{50}\) and is not reported as a
domain-shift result.

\begin{table}[!t]
\centering
\caption{Generalization under monitoring domain shifts. Not executed in
this public-data run.}
\label{tab:domain}
\begin{tabular}{lcccc}
\hline
Train \(\rightarrow\) Test & Real AP\(_{50}\) & Real Rec. & Mixed AP\(_{50}\) & Mixed Rec. \\
\hline
Summer \(\rightarrow\) Winter & --- & --- & --- & --- \\
Day \(\rightarrow\) Dusk & --- & --- & --- & --- \\
Clear \(\rightarrow\) Haze/Fog & --- & --- & --- & --- \\
\hline
\end{tabular}
\end{table}

\subsection{Qualitative Comparison and Ablations}
\label{sec:results-qual}

Fig.~\ref{fig:qual} compares FireSmokeGenNet with alpha-blend compositing
under matched backgrounds, masks, prompts, and DDIM settings. At 64\,px
the plumes occupy small regions, consistent with the Pyro-SDIS box-area
distribution. Where smoke is visible, the generator produces a different
interior appearance than constant-color blending rather than a copy of the
background.

\begin{figure*}[!t]
\centering
\includegraphics[width=0.95\textwidth]{qualitative.png}
\caption{Qualitative comparison under matched generation conditions. Each
row uses the same smoke-free background and target mask. Columns:
background, alpha-blend compositing, FireSmokeGenNet.}
\label{fig:qual}
\end{figure*}

JCA-off, MRDL-off, random encoder, unfiltered pool, and the full model
were trained under the compact detector protocol. All five configurations
produced \(\mathrm{AP}_{50}=0\). Table~\ref{tab:ablation} is therefore
unfilled. The modules remain implemented; they cannot be ranked by AP on
this host. The MRDL-weight sweep \(\omega\in\{0.0,0.4,1.0\}\) likewise
yielded \(\mathrm{AP}_{50}=0\) at every setting, so no AP curve is
reported. The training loss still uses
\((1-\omega)\mathcal{L}_{\mathrm{diff}}+\omega\mathcal{L}_{\mathrm{MRDL}}\)
with stop-gradient on the morphologically perturbed branch.

\begin{table}[!t]
\centering
\caption{Fixed-seed ablation of architectural components. AP columns are
left blank because the compact detector \(\mathrm{AP}_{50}\) was
degenerate.}
\label{tab:ablation}
\begin{tabular}{lcccc}
\hline
Configuration & AP\(_{50}\) & AP\(_{50:95}\) & Prec. & Rec. \\
\hline
No JCA (CA only) & --- & --- & --- & --- \\
No MRDL (\(\omega=0\)) & --- & --- & --- & --- \\
Encoder without ImageNet pretraining & --- & --- & --- & --- \\
Unfiltered candidate pool & --- & --- & --- & --- \\
Full FireSmokeGenNet & --- & --- & --- & --- \\
\hline
\end{tabular}
\end{table}

Table~\ref{tab:filter} reports candidate-pool sizes on the compact
percentage grid. AP and recall are unfilled. The 30\% operating point used
for generation retained 29 of 96 candidates (the configured 0.30 fraction
on the sampled pool).

\begin{table}[!t]
\centering
\caption{Retention-threshold sweep. Dataset size is measured; AP and
recall are not reported.}
\label{tab:filter}
\begin{tabular}{lccc}
\hline
Retention threshold & Dataset size & AP\(_{50}\) & Recall \\
\hline
10\% & 81 & --- & --- \\
20\% & 90 & --- & --- \\
30\% & 100 & --- & --- \\
40\% & 109 & --- & --- \\
50\% & 119 & --- & --- \\
No filtering & 167 & --- & --- \\
\hline
\end{tabular}
\end{table}

\subsection{Discussion}
\label{sec:results-disc}

\subsubsection{What this run supports}
By processing the mask and masked image through separate encoder
branches, dual-branch generation is closer to real Pyro-SDIS smoke than
alpha-blend compositing in Feature-FD and MMD\(^2\)
(Table~\ref{tab:dist}). Boundary softness is nearer the real reference
than compositing (Table~\ref{tab:bound}), while KL remains comparable at
this capacity. The three-axis composite and top-30\% keep rule were
executed; student--teacher \(\rho\) is too weak to claim a calibrated VLM
(Table~\ref{tab:vlm}). Mixed training raised five-seed classification
accuracy by \(+6.98\) percentage points (Table~\ref{tab:downstream}). The
paired test is not significant at \(\alpha=0.05\) on \(n=5\). This is a
directional proxy, not a YOLO \(\mathrm{AP}_{50}\) result.

\subsubsection{What this run does not support}
The compact execution does not support paper-scale background PSNR/SSIM
(\(28.10\)\,dB / \(0.88\)), YOLOv6--v13 \(\mathrm{AP}_{50}\) (including
the mixed YOLOv13 value of \(80.88\%\)), domain-shift AP on registered
wildfire-camera seasons, component / \(\omega\) / retention-threshold
rankings by AP, or Qwen2-VL-7B MAE/Spearman from the original VLM table.

\subsubsection{Limitations specific to the compact execution}
Tiny smoke boxes make IoU-0.5 detection degenerate. 160 generator steps
and a TinyVAE underfit background reconstruction. Heuristic quality
labels are not human ratings. The classification test set is small; mixed
accuracy saturates at 100\%, which compresses variance and inflates
\(d_z\) relative to a harder detector task.

\subsubsection{How to obtain the missing tables}
Run the paper-scale configuration on an A100-class GPU with official YOLO
weights, CLIP ViT-H/14, Qwen2-VL-7B, and the registered FLAME / HPWREN /
SMOKE5K sources. Until that run exists, those cells stay empty.
 after setting \graphicspath{{paper/}}.
%
% SOURCE OF NUMBERS: results/json/ from
%   PYTHONPATH=src python3 -u scripts/run_pipeline.py --config configs/compact.yaml
% Regenerated tables: python scripts/export_results_tables.py
%
% This fragment REPLACES the previous Results section. Methods that were not
% executed in the compact public-data run are omitted (not copied from an
% earlier draft).
% =============================================================================

\section{Results}
\label{sec:results}

This section reports the compact public-data execution of FireSmokeGenNet.
It replaces the previous manuscript Results tables. All numerical values
are taken from the JSON measurements written by the compact pipeline.
Methods that were not run---Stable Diffusion, PowerPaint, BLD,
FlameDiffuser, official YOLO v6--v13 at \(640\times 640\), and registered
FLAME / HPWREN / SMOKE5K domain splits---are omitted rather than copied
from an earlier draft.

Images were \(64\times 64\) with a stride-4 TinyVAE (latent \(16\times 16\)).
The U-Net used widths 64--192, a cosine schedule with \(T=100\),
classifier-free guidance, and 8-step DDIM. The generator was trained for
160 iterations with MRDL \(\omega=0.4\). Quality ranking used the paper
weights \(Q=0.4\,\mathrm{color}+0.4\,\mathrm{visibility}+0.2\,\mathrm{translucency}\)
and retained the top 30\% of candidates. Source-level 80/10/10 splits used
camera identifiers as the source key, with SHA-256 exact-duplicate removal
and perceptual-hash near-duplicate control. Downstream inference used five
matched seeds \(\{42,123,3407,2025,2026\}\). Public substitutes were
Pyro-SDIS smoke frames and Wikimedia landscape backgrounds. Masks were
obtained by GrabCut inside the YOLO box followed by the largest connected
component. The host was CPU-only. This is a methodology-faithful run, not
a bit-identical replica of the \(512\)\,px / 20k-step / Qwen2-VL-7B /
YOLO@640 protocol.

Table~\ref{tab:protocol} records the realized sample counts.

\begin{table}[!t]
\centering
\caption{Compact public-data protocol actually executed.}
\label{tab:protocol}
\begin{tabular}{lr}
\hline
Quantity & Value \\
\hline
Real smoke train / val / test & 74 / 10 / 23 \\
Eligible backgrounds & 90 \\
Available masks & 74 \\
Generation pairs & 48 \\
Candidates & 96 \\
Retained (top 30\%) & 29 \\
\hline
\end{tabular}
\end{table}

\subsection{Vision--Language Quality-Predictor Validation}
\label{sec:results-vlm}

The original protocol specifies a LoRA-tuned Qwen2-VL-7B ranker on 150
human-rated images. That 7B teacher does not fit in the present CPU memory
budget. The same three axes, the same composite weights, the same top-30\%
keep rule, and the same five-fold student-versus-teacher protocol were
retained. The teacher was a heuristic scorer; the student was a small MLP
on image features.

Table~\ref{tab:vlm} reports mean \(\pm\) sample standard deviation across
the five folds. The student follows the teacher most closely on visibility
(MAE \(0.97\) on a 0--10 scale) and least closely on color (MAE \(2.71\)).
Composite rank correlation is weak and high-variance
(\(\rho=0.11\pm 0.25\); fold-wise
\(\rho\in\{0.45,-0.25,0.06,0.20,0.06\}\)).
These measurements support the pipeline structure---score three axes, form
\(Q\), keep the top 30\%---but do not reproduce the original Qwen2-VL-7B
agreement (\(\rho=0.83\)). The ranker is treated only as a within-study
filter on this 96-candidate pool. Fig.~\ref{fig:vlm} visualizes per-axis MAE.

\begin{table}[!t]
\centering
\caption{Five-fold cross-validation of the task-specific ranker.
Values are mean \(\pm\) sample SD across folds.
Spearman \(\rho\) is computed on the composite score \(Q\).}
\label{tab:vlm}
\begin{tabular}{lccc}
\hline
Quality dimension & MAE \(\downarrow\) & RMSE \(\downarrow\) & Spearman \(\rho\) \(\uparrow\) \\
\hline
Color fidelity & \(2.71\pm 0.47\) & \(3.64\pm 0.48\) & --- \\
Visibility & \(0.97\pm 0.18\) & \(1.26\pm 0.28\) & --- \\
Translucency & \(1.69\pm 0.26\) & \(2.16\pm 0.22\) & --- \\
Composite score & --- & --- & \(0.11\pm 0.25\) \\
\hline
\end{tabular}
\end{table}

\begin{figure}[!t]
\centering
\includegraphics[width=0.48\textwidth]{vlm_validation.png}
\caption{Five-fold MAE of the MLP student versus the three-axis heuristic
teacher. Error bars denote one sample standard deviation across folds.}
\label{fig:vlm}
\end{figure}

\subsection{Background Preservation and Text--Image Alignment}
\label{sec:results-bg}

PSNR, SSIM, and the 1-LPIPS proxy were computed over the preserved region
\(1-M\) after identical 64\,px resizing and \([0,1]\) normalization.
Prompt similarity is cosine similarity between the generated RGB image and
a bag-of-words embedding of the caption (CLIP ViT-H/14 was not loaded).
Only methods that were actually sampled are reported. Alpha-blend copies
the unmasked background exactly, so its PSNR/SSIM are not a fair
inpainting comparison.

\begin{table}[!t]
\centering
\caption{Background preservation and text--image alignment on the common
compact manifest. PSNR, SSIM, and the 1-LPIPS proxy are computed over
\(1-M\); prompt-sim is computed over the complete image.}
\label{tab:bg}
\begin{tabular}{lcccc}
\hline
Method & PSNR \(\uparrow\) & SSIM \(\uparrow\) & 1-LPIPS \(\uparrow\) & Prompt-sim \(\uparrow\) \\
\hline
Alpha-blend compositing & 60.00 & 1.00 & --- & --- \\
FireSmokeGenNet (raw decode) & 10.08 & 0.15 & 0.80 & 0.34 \\
\hline
\end{tabular}
\end{table}

The dual-branch generator, trained for 160 steps from a randomly
initialized TinyVAE, does not preserve the background at the original
paper-scale fidelity (PSNR \(28.10\)\,dB, SSIM \(0.88\)). The
\(10.08\)\,dB / \(0.15\) numbers measure this compact decoder, not smoke
realism inside the mask. They are reported because the metric definitions
match the experimental protocol; they are not a claim that the 64\,px
model matches SD-2 inpainting.

\subsection{Distribution Alignment}
\label{sec:results-dist}

Table~\ref{tab:dist} reports Fr\'echet distance and linear MMD\(^2\) in a
shared \(16\times 16\) RGB embedding of real test smoke versus synthetic
smoke. This is the compact analogue of CLIP ViT-H/14 CLIP-FD / linear MMD.
Lower is better. GAN, SD-inpainting, and FlameDiffuser rows are omitted
because those generators were not run.

\begin{table}[!t]
\centering
\caption{Feature-space distance between real Pyro-SDIS test smoke and
synthetic smoke. Lower is better. Features are \(16\times 16\) RGB, not
CLIP ViT-H/14.}
\label{tab:dist}
\begin{tabular}{lcc}
\hline
Method & Feature-FD \(\downarrow\) & MMD\(^2\) \(\downarrow\) \\
\hline
Alpha-blend compositing & 48.55 & 33.28 \\
FireSmokeGenNet & \textbf{44.54} & \textbf{28.83} \\
\hline
\end{tabular}
\end{table}

FireSmokeGenNet sits closer to real smoke than training-free compositing
on both distances (\(\Delta\) Feature-FD \(=-4.01\),
\(\Delta\) MMD\(^2=-4.45\)). The comparison is a single-manifest
descriptive result with no image-level confidence interval.
Fig.~\ref{fig:dist} visualizes the same values.

\begin{figure}[!t]
\centering
\includegraphics[width=0.48\textwidth]{distribution.png}
\caption{Feature-FD and linear MMD\(^2\) between real test smoke and
synthetic smoke. Lower is better.}
\label{fig:dist}
\end{figure}

\subsection{Boundary Softness Analysis}
\label{sec:results-bound}

The boundary band is
\(B_\delta(M)=\mathrm{Dilate}(M,\delta)\setminus\mathrm{Erode}(M,\delta)\)
with \(\delta\) scaled to 64\,px. Softness \(S\) is the mean gradient
magnitude in that band. KL uses 100-bin histograms of \(\|\nabla I\|_2\)
with a \(10^{-8}\) floor.

\begin{table}[!t]
\centering
\caption{Boundary-quality metrics. Real-smoke softness is a reference, not
a method.}
\label{tab:bound}
\begin{tabular}{lcc}
\hline
Method & Boundary softness \(S\) & KL\((\mathrm{real}\,\|\,\mathrm{synth})\) \(\downarrow\) \\
\hline
Real smoke & 0.363 & --- \\
Alpha-blend compositing & 0.438 & 15.20 \\
FireSmokeGenNet & \textbf{0.310} & 15.31 \\
\hline
\end{tabular}
\end{table}

FireSmokeGenNet is nearer the real reference than compositing
(\(0.310\) versus \(0.363\); blend \(0.438\)), i.e., less abrupt
mask-edge gradients. KL is comparable at this capacity (\(15.31\) versus
\(15.20\)). The originally published gap (softness \(0.124\) versus real
\(0.119\); KL \(0.041\) versus \(0.187\)) is not recovered at 64\,px /
160 steps. Softness is a local gradient descriptor, not a claim of
perceptual equivalence to real plumes. Fig.~\ref{fig:bound} shows the
three softness values.

\begin{figure}[!t]
\centering
\includegraphics[width=0.48\textwidth]{boundary.png}
\caption{Boundary-gradient magnitude in the \(\delta\)-band. Real smoke is
a reference distribution, not a generation method.}
\label{fig:bound}
\end{figure}

\subsection{Downstream Detection Performance}
\label{sec:results-det}

Official YOLO families at 640\,px were not executed. Pyro-SDIS median box
area is approximately \(0.17\%\) of the image; at 64--128\,px the box is
on the order of 10\,px, and IoU \(0.5\) is unstable. All eight compact
family stand-ins produced \(\mathrm{AP}_{50}=0\). Those zeros are not
filled into Table~\ref{tab:downstream}, and the previously published
YOLOv13 numbers (\(78.16\pm 0.27\%\rightarrow 80.88\pm 0.70\%\)) are not
copied.

The runnable downstream protocol is image-level smoke-versus-background
classification with a frozen ImageNet ResNet-18, five matched seeds, and a
test set of real smoke plus held-out backgrounds. Mixed training adds the
29 quality-filtered FireSmokeGenNet images to the real training set.

\begin{table*}[!t]
\centering
\caption{Five-seed comparison of real-only and mixed-data training.
Performance is smoke/background accuracy (\%). The upper endpoint of the
real-only \(t\)-interval exceeds 100\% because \(n=5\) and the sample SD
is large; the feasible range is \([82.22, 100]\).}
\label{tab:downstream}
\begin{tabular}{lccccc ccc ccc}
\hline
Condition & s1 & s2 & s3 & s4 & s5 & Mean \(\pm\) SD & 95\% CI & \(\Delta\) & Test & \(p\) & \(d_z\) \\
\hline
Real only
  & 81.40 & 86.05 & 97.67 & 100.00 & 100.00
  & \(93.02\pm 8.70\) & [82.22, 103.83]
  & --- & --- & --- & --- \\
Mixed data
  & 100.00 & 100.00 & 100.00 & 100.00 & 100.00
  & \(100.00\pm 0.00\) & [100.00, 100.00]
  & \(+6.98\) & Paired \(t\) & 0.147 & 0.80 \\
\hline
\end{tabular}
\end{table*}

The mixed-data mean is higher by 6.98 percentage points. Shapiro--Wilk on
the paired differences did not reject normality (\(p=0.111\)), so a
two-sided paired \(t\)-test was used. The test is not significant at
\(\alpha=0.05\) (\(p=0.147\)). Cohen's \(d_z=0.80\) is a large
standardized paired effect with only five runs; it is reported together
with the \(p\)-value and interval, not as a substitute for a detector AP
table. Holm correction is not applied because only one downstream
comparison was valid. Fig.~\ref{fig:cls} visualizes the two five-seed
means.

\begin{figure}[!t]
\centering
\includegraphics[width=0.48\textwidth]{classification.png}
\caption{Five-seed image-level smoke/background accuracy. Error bars
denote one sample standard deviation. Mixed training adds 29
quality-filtered FireSmokeGenNet images.}
\label{fig:cls}
\end{figure}

\subsection{Domain Generalization Under Environmental Shifts}
\label{sec:results-domain}

Table~\ref{tab:domain} is left empty. Summer\(\rightarrow\)winter,
day\(\rightarrow\)dusk, and clear\(\rightarrow\)haze splits on FLAME /
HPWREN / SMOKE5K were not downloaded because those sources require
registration. A compact clear-to-haze AP probe on the public subset also
returned degenerate \(\mathrm{AP}_{50}\) and is not reported as a
domain-shift result.

\begin{table}[!t]
\centering
\caption{Generalization under monitoring domain shifts. Not executed in
this public-data run.}
\label{tab:domain}
\begin{tabular}{lcccc}
\hline
Train \(\rightarrow\) Test & Real AP\(_{50}\) & Real Rec. & Mixed AP\(_{50}\) & Mixed Rec. \\
\hline
Summer \(\rightarrow\) Winter & --- & --- & --- & --- \\
Day \(\rightarrow\) Dusk & --- & --- & --- & --- \\
Clear \(\rightarrow\) Haze/Fog & --- & --- & --- & --- \\
\hline
\end{tabular}
\end{table}

\subsection{Qualitative Comparison and Ablations}
\label{sec:results-qual}

Fig.~\ref{fig:qual} compares FireSmokeGenNet with alpha-blend compositing
under matched backgrounds, masks, prompts, and DDIM settings. At 64\,px
the plumes occupy small regions, consistent with the Pyro-SDIS box-area
distribution. Where smoke is visible, the generator produces a different
interior appearance than constant-color blending rather than a copy of the
background.

\begin{figure*}[!t]
\centering
\includegraphics[width=0.95\textwidth]{qualitative.png}
\caption{Qualitative comparison under matched generation conditions. Each
row uses the same smoke-free background and target mask. Columns:
background, alpha-blend compositing, FireSmokeGenNet.}
\label{fig:qual}
\end{figure*}

JCA-off, MRDL-off, random encoder, unfiltered pool, and the full model
were trained under the compact detector protocol. All five configurations
produced \(\mathrm{AP}_{50}=0\). Table~\ref{tab:ablation} is therefore
unfilled. The modules remain implemented; they cannot be ranked by AP on
this host. The MRDL-weight sweep \(\omega\in\{0.0,0.4,1.0\}\) likewise
yielded \(\mathrm{AP}_{50}=0\) at every setting, so no AP curve is
reported. The training loss still uses
\((1-\omega)\mathcal{L}_{\mathrm{diff}}+\omega\mathcal{L}_{\mathrm{MRDL}}\)
with stop-gradient on the morphologically perturbed branch.

\begin{table}[!t]
\centering
\caption{Fixed-seed ablation of architectural components. AP columns are
left blank because the compact detector \(\mathrm{AP}_{50}\) was
degenerate.}
\label{tab:ablation}
\begin{tabular}{lcccc}
\hline
Configuration & AP\(_{50}\) & AP\(_{50:95}\) & Prec. & Rec. \\
\hline
No JCA (CA only) & --- & --- & --- & --- \\
No MRDL (\(\omega=0\)) & --- & --- & --- & --- \\
Encoder without ImageNet pretraining & --- & --- & --- & --- \\
Unfiltered candidate pool & --- & --- & --- & --- \\
Full FireSmokeGenNet & --- & --- & --- & --- \\
\hline
\end{tabular}
\end{table}

Table~\ref{tab:filter} reports candidate-pool sizes on the compact
percentage grid. AP and recall are unfilled. The 30\% operating point used
for generation retained 29 of 96 candidates (the configured 0.30 fraction
on the sampled pool).

\begin{table}[!t]
\centering
\caption{Retention-threshold sweep. Dataset size is measured; AP and
recall are not reported.}
\label{tab:filter}
\begin{tabular}{lccc}
\hline
Retention threshold & Dataset size & AP\(_{50}\) & Recall \\
\hline
10\% & 81 & --- & --- \\
20\% & 90 & --- & --- \\
30\% & 100 & --- & --- \\
40\% & 109 & --- & --- \\
50\% & 119 & --- & --- \\
No filtering & 167 & --- & --- \\
\hline
\end{tabular}
\end{table}

\subsection{Discussion}
\label{sec:results-disc}

\subsubsection{What this run supports}
By processing the mask and masked image through separate encoder
branches, dual-branch generation is closer to real Pyro-SDIS smoke than
alpha-blend compositing in Feature-FD and MMD\(^2\)
(Table~\ref{tab:dist}). Boundary softness is nearer the real reference
than compositing (Table~\ref{tab:bound}), while KL remains comparable at
this capacity. The three-axis composite and top-30\% keep rule were
executed; student--teacher \(\rho\) is too weak to claim a calibrated VLM
(Table~\ref{tab:vlm}). Mixed training raised five-seed classification
accuracy by \(+6.98\) percentage points (Table~\ref{tab:downstream}). The
paired test is not significant at \(\alpha=0.05\) on \(n=5\). This is a
directional proxy, not a YOLO \(\mathrm{AP}_{50}\) result.

\subsubsection{What this run does not support}
The compact execution does not support paper-scale background PSNR/SSIM
(\(28.10\)\,dB / \(0.88\)), YOLOv6--v13 \(\mathrm{AP}_{50}\) (including
the mixed YOLOv13 value of \(80.88\%\)), domain-shift AP on registered
wildfire-camera seasons, component / \(\omega\) / retention-threshold
rankings by AP, or Qwen2-VL-7B MAE/Spearman from the original VLM table.

\subsubsection{Limitations specific to the compact execution}
Tiny smoke boxes make IoU-0.5 detection degenerate. 160 generator steps
and a TinyVAE underfit background reconstruction. Heuristic quality
labels are not human ratings. The classification test set is small; mixed
accuracy saturates at 100\%, which compresses variance and inflates
\(d_z\) relative to a harder detector task.

\subsubsection{How to obtain the missing tables}
Run the paper-scale configuration on an A100-class GPU with official YOLO
weights, CLIP ViT-H/14, Qwen2-VL-7B, and the registered FLAME / HPWREN /
SMOKE5K sources. Until that run exists, those cells stay empty.
